\scanpage{172}\mspage{117}
associée à $X$. On peut oublier le scindage, et même (pour \unsure{calculer|étudier} les $\mathrm{Hom}_{\mathrm{HOT}_W(S)}(X,Y)$) travailler avec des $S$-foncteurs \uline{quelconques} entre catégories fibrées sur $S$ -- pas la peine même de les supposer cartésiens, voire $=$ compatibles avec des scindages. Cela me paraît une chose \unsure{fort|très} \unsure{importante} pour le formalisme des $\mathrm{HOT}_W(S)$, et très spéciale au cas \deleted{\unsure{des}} de la catégorie de modèles $\mathrm{Cat}$.
C'est \inserted{un peu} \unsure{tacite|traité} dans le théorème-scholie -- mais mes doutes initiaux concernant les catégories $\underline{\mathrm{Fibsc}}$ elles-mêmes montrent bien que je sentais qu'il y avait là quelque chose de pas évident, ni même \unsure{(pour moi)|/ pas gros|/ par gros} de vraisemblable. Et pourtant l'existence du foncteur $\unsure{\Phi|Q}$ montre qu'il en est bien ainsi.

On peut donc interpréter $f^*$ en termes de $\underline{\mathrm{Fib}}\,S$, $\underline{\mathrm{Fib}}\,S'$, par le foncteur chg$^{\mathrm t}$ de base
\[
f^*:\underline{\mathrm{Fib}}\,S\longrightarrow\underline{\mathrm{Fib}}\,S',\qquad X\longmapsto X\times_S S'\,/S',
\]
lequel foncteur préserve évidemment les $W$-équivalences par fibres $W_S^f$, $W_{S'}^f$, qui est celle \unsure{pertinentes} ici
\[
f^*(W_S^f)\subset W_{S'}^f\deleted{\unsure{\}}},\quad\unsure{d}
\]
de sorte que $f^*$ induit bien
\[
(\underline{\mathrm{Fib}}\,S)(W_S^f)^{-1}\longrightarrow(\underline{\mathrm{Fib}}\,S')(W_S^f)^{-1},
\]
% archive sheet number: 171
